\section{Problem Set 8: Sequences and Functions}

\begin{enumerate}
\item \textbf{Problem 1}\par
Decide which of the following assignments define functions:
\begin{enumerate}
\item each real number \(x\) is assigned \(x^2\);
\item each positive number \(x\) is assigned both \(\sqrt{x}\) and \(-\sqrt{x}\);
\item each student in a class is assigned the student's date of birth.
\end{enumerate}
Explain each decision.

\item \textbf{Problem 2}\par
For
\[
f:\mathbb{R}\to\mathbb{R},
\qquad
f(x)=x^2-4,
\]
state the domain and codomain, and determine the range.

\item \textbf{Problem 3}\par
Let
\[
f(x)=2x^2-3x+1.
\]
Compute \(f(0)\), \(f(-2)\), \(f(a)\), and \(f(x+h)\).

\item \textbf{Problem 4}\par
If
\[
g(x)=\frac{1}{x-3},
\]
determine the natural real domain of \(g\) and compute \(g(1)\).

\item \textbf{Problem 5}\par
The graph of a function contains the points
\[
(-2,3),\quad (0,-1),\quad (4,5).
\]
State the corresponding function values.

\item \textbf{Problem 6}\par
For
\[
f(x)=x^2-5x+6,
\]
find the \(y\)-intercept and all \(x\)-intercepts.

\item \textbf{Problem 7}\par
Determine the natural real domain of each function:
\[
f(x)=\sqrt{x+2},
\qquad
g(x)=\frac{x+1}{x^2-9},
\qquad
h(x)=\ln(4-x).
\]

\item \textbf{Problem 8}\par
Find the range of each function:
\[
f(x)=x^2+2,
\qquad
g(x)=-3x+1,
\qquad
h(x)=\sqrt{x}.
\]

\item \textbf{Problem 9}\par
Classify each function as constant, affine, quadratic, polynomial, rational, exponential, logarithmic, or trigonometric:
\[
3,\quad 2x-5,\quad x^2+1,\quad \frac{x+1}{x-2},\quad 4^x,\quad \ln x,\quad \sin x.
\]

\item \textbf{Problem 10}\par
Determine which of the functions
\[
f(x)=3x,\qquad g(x)=3x+2,\qquad h(x)=-x
\]
are linear in the sense \(f(x+y)=f(x)+f(y)\). Explain why an affine function need not be linear.

\item \textbf{Problem 11}\par
For
\[
f(x)=2^x
\quad\text{and}\quad
g(x)=\log_2 x,
\]
compute \(f(-2)\), \(f(3)\), \(g(1)\), \(g(8)\), and explain the relationship between the two functions.

\item \textbf{Problem 12}\par
Let
\[
a_n=3n-1,\qquad n\ge1.
\]
Write the first five terms and compute \(a_{20}\).

\item \textbf{Problem 13}\par
A sequence is defined by
\[
a_1=4,\qquad a_{n+1}=a_n-3.
\]
Write the first six terms.

\item \textbf{Problem 14}\par
Find an explicit formula for the arithmetic sequence
\[
5,\ 9,\ 13,\ 17,\ldots
\]

\item \textbf{Problem 15}\par
Find an explicit formula for the geometric sequence
\[
3,\ 6,\ 12,\ 24,\ldots
\]

\item \textbf{Problem 16}\par
Define recursively the sequence
\[
2,\ 6,\ 18,\ 54,\ldots
\]
and state the starting index.

\item \textbf{Problem 17}\par
For
\[
a_n=\frac{1}{n},
\]
list the first five graph points \((n,a_n)\). Explain why the graph should normally be drawn as separated points.

\item \textbf{Problem 18}\par
Compare the domains of
\[
a_n=n^2,\qquad n\in\mathbb{N},
\]
and
\[
f(x)=x^2,\qquad x\in\mathbb{R}.
\]
Which expressions \(a_{3/2}\) and \(f(3/2)\) are meaningful?

\item \textbf{Problem 19}\par
A graph rises on \((-3,1)\), falls on \((1,4)\), crosses the \(x\)-axis at \(-2\) and \(3\), and has a local maximum at \((1,5)\). Translate this information into statements about the function.

\item \textbf{Problem 20}\par
A taxi fare consists of a fixed fee of \(6\) units and \(2.5\) units for every kilometre travelled. Define a function \(C(d)\), state a meaningful domain, compute \(C(8)\), and explain whether the function is linear or affine.
\end{enumerate}
